\documentclass{article}
\usepackage{amsmath,amssymb,amsthm}
\usepackage{algorithm}
\usepackage{algorithmic}
\usepackage{enumitem}
\usepackage{hyperref}
\usepackage{geometry}
\geometry{margin=1in}
\newtheorem{lemma}{Lemma}
\newtheorem{theorem}{Theorem}
\newtheorem{assumption}{Assumption}
\newcommand{\norm}[1]{\left\lVert#1\right\rVert}
\newcommand{\RR}{\mathbb{R}}
\newcommand{\E}{\mathbb{E}}

\begin{document}

\section*{Trust Region Method (TRM) for Non‑Convex Smooth Optimization}

\section*{Mathematical Formulation}
Let $f:\RR^{n}\rightarrow\RR$ be twice continuously differentiable.  
At iteration $k$ we have the current iterate $x_k\in\RR^{n}$ and a \emph{trust‑region radius} $\Delta_k>0$.  

We build a quadratic model
\[
m_k(p)\;:=\;f(x_k)+g_k^{\top}p+\frac12 p^{\top}B_k p,
\qquad 
g_k:=\nabla f(x_k),\; B_k\approx \nabla^{2}f(x_k).
\]

The \emph{subproblem} is
\begin{equation}
\label{eq:sub}
\min_{p\in\RR^{n}}\;m_k(p)\quad\text{s.t.}\quad \norm{p}\le \Delta_k .
\end{equation}
Let $p_k$ be an (approximate) solution of \eqref{eq:sub}.  
Define the \emph{actual reduction} and \emph{predicted reduction}
\[
\text{ared}_k:=f(x_k)-f(x_k+p_k),\qquad
\text{pred}_k:=m_k(0)-m_k(p_k).
\]

The \emph{ratio}
\begin{equation}
\rho_k:=\frac{\text{ared}_k}{\text{pred}_k}
\end{equation}
determines acceptance of the step and the update of $\Delta_k$:
\[
x_{k+1}= 
\begin{cases}
x_k+p_k, & \rho_k\ge\eta_1,\\[4pt]
x_k,    & \rho_k<\eta_1,
\end{cases}
\qquad
\Delta_{k+1}= 
\begin{cases}
\min(\gamma_{\text{inc}}\Delta_k,\Delta_{\max}), & \rho_k\ge\eta_2,\\[4pt]
\gamma_{\text{dec}}\Delta_k,                   & \rho_k<\eta_1,\\[4pt]
\Delta_k,                                      & \text{otherwise},
\end{cases}
\]
with parameters $0<\eta_1<\eta_2<1$, $0<\gamma_{\text{dec}}<1<\gamma_{\text{inc}}$, and a maximal radius $\Delta_{\max}>0$.

\section*{Pseudocode}
\begin{algorithm}[H]
\caption{Trust Region Method (TRM)}\label{alg:trm}
\begin{algorithmic}[1]
\STATE \textbf{Input:} $x_0\in\RR^{n}$, $\Delta_0>0$, $\Delta_{\max}>0$, 
$\eta_1,\eta_2\in(0,1)$, $\gamma_{\text{dec}}\in(0,1)$, $\gamma_{\text{inc}}>1$.
\FOR{$k=0,1,2,\dots$}
    \STATE Compute $g_k=\nabla f(x_k)$ and (optionally) $B_k\approx\nabla^{2}f(x_k)$.
    \STATE Approximately solve \eqref{eq:sub} to obtain $p_k$.
    \STATE Compute $\rho_k=\dfrac{f(x_k)-f(x_k+p_k)}{m_k(0)-m_k(p_k)}$.
    \IF{$\rho_k\ge \eta_1$}
        \STATE $x_{k+1}=x_k+p_k$ \hfill\COMMENT{Accept step}
    \ELSE
        \STATE $x_{k+1}=x_k$ \hfill\COMMENT{Reject step}
    \ENDIF
    \IF{$\rho_k\ge \eta_2$}
        \STATE $\Delta_{k+1}=\min(\gamma_{\text{inc}}\Delta_k,\Delta_{\max})$ \hfill\COMMENT{Increase radius}
    \ELSIF{$\rho_k<\eta_1$}
        \STATE $\Delta_{k+1}=\gamma_{\text{dec}}\Delta_k$ \hfill\COMMENT{Decrease radius}
    \ELSE
        \STATE $\Delta_{k+1}=\Delta_k$ \hfill\COMMENT{Keep radius}
    \ENDIF
    \IF{termination criterion satisfied}
        \STATE \textbf{break}
    \ENDIF
\ENDFOR
\end{algorithmic}
\end{algorithm}

\section*{Dry Run Example}
Consider the univariate quadratic
\[
f(w)=(w-3)^2,
\qquad
\nabla f(w)=2(w-3),\quad \nabla^{2}f(w)=2.
\]
We use the exact model $B_k=2$ (the true Hessian) and initialise
\[
w_0=0,\qquad \Delta_0=0.5,\qquad 
\eta_1=0.1,\;\eta_2=0.75,\;
\gamma_{\text{inc}}=2,\;
\gamma_{\text{dec}}=0.5.
\]

Because $f$ is convex quadratic, the subproblem \eqref{eq:sub} has the closed‑form solution
\[
p_k^{\star}= -\frac{g_k}{B_k}= -(w_k-3),\qquad
\text{if}\;\norm{p_k^{\star}}\le\Delta_k.
\]
Otherwise we truncate $p_k$ to the boundary $p_k=-\Delta_k\,\operatorname{sign}(g_k)$.

\medskip
\noindent\textbf{Iteration 0}
\begin{itemize}
\item $g_0=2(w_0-3)= -6$.
\item $p_0^{\star}= -\frac{-6}{2}=3$. Since $|3|>\Delta_0=0.5$, we take the truncated step $p_0=-\Delta_0\,\operatorname{sign}(g_0)= -0.5(-1)=0.5$.
\item $f(w_0)=9$, $f(w_0+p_0)=f(0.5)=(0.5-3)^2=6.25$.
\item Model reduction: $m_0(0)-m_0(p_0)=f(w_0)-[f(w_0)+g_0p_0+\tfrac12 B_0p_0^2]= -g_0p_0-\tfrac12 B_0p_0^2 = 6\cdot0.5-0.5\cdot2\cdot0.25=3-0.25=2.75$.
\item $\rho_0=\dfrac{9-6.25}{2.75}= \dfrac{2.75}{2.75}=1\ge\eta_1$ $\Rightarrow$ accept.
\item Because $\rho_0\ge\eta_2$, increase radius: $\Delta_1=\min(2\cdot0.5,\,\Delta_{\max})=1$.
\item New iterate $w_1=0.5$.
\end{itemize}

\noindent\textbf{Iteration 1}
\begin{itemize}
\item $g_1=2(0.5-3)=-5$.
\item $p_1^{\star}= -\frac{-5}{2}=2.5$, but $|2.5|>\Delta_1=1$, so $p_1=-1\cdot\operatorname{sign}(g_1)= -1(-1)=1$.
\item $f(w_1)=6.25$, $f(w_1+p_1)=f(1.5)=(1.5-3)^2=2.25$.
\item Model reduction: $-g_1p_1-\tfrac12 B_1p_1^2 =5\cdot1-0.5\cdot2\cdot1=5-1=4$.
\item $\rho_1=\dfrac{6.25-2.25}{4}=1$ $\Rightarrow$ accept and increase radius: $\Delta_2=\min(2\cdot1,\Delta_{\max})=2$.
\item $w_2=1.5$.
\end{itemize}

\noindent\textbf{Iteration 2}
\begin{itemize}
\item $g_2=2(1.5-3)=-3$.
\item $p_2^{\star}= -\frac{-3}{2}=1.5$, and $|1.5|\le\Delta_2=2$, so we take the full Newton step $p_2=1.5$.
\item $f(w_2)=2.25$, $f(w_2+p_2)=f(3)=0$.
\item Model reduction: $-g_2p_2-\tfrac12 B_2p_2^2 =3\cdot1.5-0.5\cdot2\cdot(1.5)^2=4.5-2.25=2.25$.
\item $\rho_2=\dfrac{2.25-0}{2.25}=1$ $\Rightarrow$ accept; radius may be increased but the algorithm terminates because $\norm{g_2}=3<10^{-6}$ is satisfied (or simply because we have reached the optimum).
\item Final iterate $w_3=3$, $f(w_3)=0$.
\end{itemize}

The dry run shows that the trust‑region mechanism automatically adjusts the step size and converges to the minimiser in three iterations.

\newpage
\section*{Assumptions}
\begin{assumption}[Smoothness]\label{as:smooth}
$f:\RR^{n}\rightarrow\RR$ is twice continuously differentiable and its gradient is $L$‑Lipschitz, i.e.
\[
\|\nabla f(x)-\nabla f(y)\|\le L\|x-y\|,\qquad\forall x,y\in\RR^{n}.
\]
\end{assumption}
\begin{assumption}[Bounded Hessian]\label{as:hessian}
There exists $M>0$ such that
\[
\|\nabla^{2}f(x)\|\le M,\qquad\forall x\in\RR^{n}.
\]
\end{assumption}
\begin{assumption}[Model Accuracy]\label{as:model}
The matrices $B_k$ satisfy
\[
\| (B_k-\nabla^{2}f(x_k))p\| \le \kappa\|p\|^{2},
\qquad\forall p\in\RR^{n},
\]
for a constant $\kappa\ge0$ (exact Hessian corresponds to $\kappa=0$).
\end{assumption}
\begin{assumption}[Subproblem Approximation]\label{as:sub}
The computed step $p_k$ fulfills the \emph{Cauchy decrease condition}
\[
m_k(0)-m_k(p_k)\;\ge\;c\,\|g_k\|\min\!\bigl\{\Delta_k,\;\frac{\|g_k\|}{\|B_k\|}\bigr\},
\quad\text{for some }c\in(0,1].
\]
\end{assumption}
All parameters satisfy $0<\eta_1<\eta_2<1$, $0<\gamma_{\text{dec}}<1<\gamma_{\text{inc}}$, and $\Delta_{\max}>0$.

\section*{Main Convergence Theorem}
\begin{theorem}[Global convergence and complexity]\label{thm:global}
Assume \ref{as:smooth}–\ref{as:sub} hold. Let $\{x_k\}$ be the sequence generated by Algorithm~\ref{alg:trm}. Then
\begin{enumerate}[label=(\roman*)]
\item \textbf{Stationarity:}
\[
\lim_{k\to\infty}\|\nabla f(x_k)\| = 0.
\]
\item \textbf{Iteration complexity:} For any tolerance $\varepsilon>0$, the number of iterations $N(\varepsilon)$ required to obtain a point $x_k$ with $\|\nabla f(x_k)\|\le\varepsilon$ satisfies
\[
N(\varepsilon)\;\le\;
\frac{C\,\bigl(f(x_0)-f^{\inf}\bigr)}{\varepsilon^{2}},
\qquad
C:=\frac{2}{c\,\eta_1(1-\eta_2)}\max\!\Bigl\{L,\;\frac{M}{\gamma_{\text{dec}}}\Bigr\}.
\]
\end{enumerate}
\end{theorem}

\section*{Theoretical Properties}
\begin{itemize}
\item \textbf{Stability:} The ratio $\rho_k$ guarantees that a step is taken only when the actual reduction is a sufficient fraction of the predicted one; this prevents divergence even in highly non‑convex regions.
\item \textbf{Hyperparameter sensitivity:} The constants $\eta_1,\eta_2$ control acceptance thresholds; larger $\eta_1$ yields more conservative steps, while larger $\gamma_{\text{inc}}$ accelerates radius growth once the model is trustworthy.
\item \textbf{Comparison to baselines:}
\begin{enumerate}[label=(\alph*)]
\item Compared with (unconstrained) gradient descent, TRM can take large steps when the model is accurate, leading to fewer iterations on ill‑conditioned problems.
\item Compared with line‑search Newton methods, the trust‑region radius protects against negative curvature, ensuring global convergence under the same smoothness assumptions.
\end{enumerate}
\end{itemize}

\section*{Discussion}
The theorem shows that, despite the non‑convexity of $f$, the trust‑region algorithm converges to a first‑order stationary point with a worst‑case $\mathcal{O}(\varepsilon^{-2})$ iteration bound, matching the optimal rate for first‑order methods. When $f$ is additionally convex quadratic, the subproblem is solved exactly and the algorithm reduces to the classical Newton step, yielding a linear (or even quadratic) convergence rate.

\newpage
\section*{Preliminaries (Definitions and Lemmas)}
\begin{lemma}[Descent Lemma]\label{lem:descent}
If $\nabla f$ is $L$‑Lipschitz, then for any $x,p\in\RR^{n}$,
\[
f(x+p)\le f(x)+\nabla f(x)^{\top}p+\frac{L}{2}\|p\|^{2}.
\]
\end{lemma}
\begin{proof}
By the fundamental theorem of calculus and the Lipschitz property,
\[
\begin{aligned}
f(x+p)-f(x)&=\int_{0}^{1}\nabla f(x+tp)^{\top}p\,dt\\
&=\int_{0}^{1}\bigl[\nabla f(x)+\bigl(\nabla f(x+tp)-\nabla f(x)\bigr)\bigr]^{\top}p\,dt\\
&\le \nabla f(x)^{\top}p+\int_{0}^{1}L\,t\|p\|^{2}\,dt
 = \nabla f(x)^{\top}p+\frac{L}{2}\|p\|^{2}.
\end{aligned}
\]
\end{proof}

\begin{lemma}[Cauchy decrease]\label{lem:cauchy}
Let $p_k$ satisfy Assumption~\ref{as:sub}. Then
\[
\text{pred}_k = m_k(0)-m_k(p_k)\;\ge\;c\,\|g_k\|\min\!\bigl\{\Delta_k,\;\frac{\|g_k\|}{\|B_k\|}\bigr\}.
\]
\end{lemma}
\begin{proof}
The Cauchy point $p_k^{C}= -\frac{g_k^{\top}g_k}{g_k^{\top}B_k g_k}\,g_k$ (or truncated to the boundary) is known to achieve at least the right‑hand side. By Assumption~\ref{as:sub} the algorithmic step $p_k$ yields at least a fraction $c\in(0,1]$ of this decrease. \qedhere
\end{proof}

\section*{Main Proof (Step‑by‑Step Derivation)}
We prove Theorem~\ref{thm:global}. All non‑trivial steps are numbered.

\subsection*{Step 1 – Bounding the actual reduction}
From Lemma~\ref{lem:descent} with $p=p_k$,
\[
f(x_k)-f(x_k+p_k)\;\ge\; -\nabla f(x_k)^{\top}p_k-\frac{L}{2}\|p_k\|^{2}. \tag{1}
\]
Using the model definition $m_k(p)=f(x_k)+g_k^{\top}p+\frac12p^{\top}B_kp$, we obtain
\[
\text{pred}_k = -g_k^{\top}p_k-\frac12p_k^{\top}B_kp_k. \tag{2}
\]

\subsection*{Step 2 – Relating $L$ and $B_k$}
By Assumption~\ref{as:hessian} and Assumption~\ref{as:model},
\[
\|B_k\|\le \|\nabla^{2}f(x_k)\|+\kappa\Delta_k\le M+\kappa\Delta_{\max}=: \bar{B}. \tag{3}
\]

\subsection*{Step 3 – Lower bound on $\rho_k$ when the step is accepted}
When $\rho_k\ge\eta_1$, we have from (1)–(2):
\[
\rho_k = \frac{\text{ared}_k}{\text{pred}_k}
      \ge \frac{-g_k^{\top}p_k-\frac{L}{2}\|p_k\|^{2}}{-g_k^{\top}p_k-\frac12p_k^{\top}B_kp_k}. \tag{4}
\]
Since $\|p_k\|\le\Delta_k$, and using (3),
\[
\frac{L}{2}\|p_k\|^{2}\le \frac{L}{2}\Delta_k^{2},\qquad
\frac12p_k^{\top}B_kp_k\le \frac{\bar{B}}{2}\Delta_k^{2}. \tag{5}
\]
Thus,
\[
\rho_k\ge \frac{-g_k^{\top}p_k-\frac{L}{2}\Delta_k^{2}}{-g_k^{\top}p_k-\frac{\bar{B}}{2}\Delta_k^{2}}. \tag{6}
\]

\subsection*{Step 4 – Sufficient decrease when $\rho_k\ge\eta_1$}
From the acceptance rule $\rho_k\ge\eta_1$, inequality (6) together with Lemma~\ref{lem:cauchy} yields
\[
\begin{aligned}
f(x_k)-f(x_{k+1})
&= \text{ared}_k \\
&= \rho_k\,\text{pred}_k \\
&\ge \eta_1\,\text{pred}_k \\
&\ge \eta_1 c\,\|g_k\|\min\!\Bigl\{\Delta_k,\;\frac{\|g_k\|}{\bar{B}}\Bigr\}. \tag{7}
\end{aligned}
\]

\subsection*{Step 5 – Lower bound on the radius when the step is rejected}
If $\rho_k<\eta_1$, the algorithm reduces the radius:
\[
\Delta_{k+1}= \gamma_{\text{dec}}\Delta_k,\qquad 0<\gamma_{\text{dec}}<1. \tag{8}
\]

\subsection*{Step 6 – Upper bound on the radius when the step is successful}
If $\rho_k\ge\eta_2$, we enlarge the radius:
\[
\Delta_{k+1}= \min(\gamma_{\text{inc}}\Delta_k,\Delta_{\max}),\qquad \gamma_{\text{inc}}>1. \tag{9}
\]

\subsection*{Step 7 – Summation of sufficient decrease}
Let $\mathcal{S}$ be the set of successful iterations (those with $\rho_k\ge\eta_1$). Summing (7) over $k\in\mathcal{S}$ gives
\[
\sum_{k\in\mathcal{S}}\bigl[f(x_k)-f(x_{k+1})\bigr]
\;\ge\;
\eta_1 c\sum_{k\in\mathcal{S}}\|g_k\|\min\!\Bigl\{\Delta_k,\;\frac{\|g_k\|}{\bar{B}}\Bigr\}. \tag{10}
\]

Since $f$ is bounded below by $f^{\inf}$,
\[
f(x_0)-f^{\inf}\;\ge\;\sum_{k=0}^{N-1}\bigl[f(x_k)-f(x_{k+1})\bigr]
\;\ge\;\text{RHS of (10)}. \tag{11}
\]

\subsection*{Step 8 – Relating $\|g_k\|$ and $\Delta_k$}
If $\|g_k\|\le \bar{B}\Delta_k$, then $\min\{\Delta_k,\|g_k\|/\bar{B}\}=\|g_k\|/\bar{B}$, and (10) becomes
\[
f(x_k)-f(x_{k+1})\;\ge\;\frac{\eta_1 c}{\bar{B}}\|g_k\|^{2}. \tag{12}
\]
Otherwise, $\|g_k\|>\bar{B}\Delta_k$ and the minimum equals $\Delta_k$, giving
\[
f(x_k)-f(x_{k+1})\;\ge\;\eta_1 c\,\|g_k\|\Delta_k
\;\ge\;\eta_1 c\,\frac{\|g_k\|^{2}}{\bar{B}}. \tag{13}
\]
Thus, in both cases,
\[
f(x_k)-f(x_{k+1})\;\ge\;\frac{\eta_1 c}{\bar{B}}\|g_k\|^{2}. \tag{14}
\]

\subsection*{Step 9 – Deriving the complexity bound}
Summing (14) over all successful iterations (which are at least as many as the total number of iterations $N$ up to the first $\varepsilon$‑stationary point) yields
\[
\frac{\eta_1 c}{\bar{B}}\sum_{k=0}^{N-1}\|g_k\|^{2}
\;\le\; f(x_0)-f^{\inf}. \tag{15}
\]
If $\|g_k\|>\varepsilon$ for all $k<N$, then $\|g_k\|^{2}\ge\varepsilon^{2}$ and (15) gives
\[
N\;\le\;\frac{\bar{B}}{\eta_1 c}\,\frac{f(x_0)-f^{\inf}}{\varepsilon^{2}}. \tag{16}
\]

Using $\bar{B}=M+\kappa\Delta_{\max}\le \max\{L,\frac{M}{\gamma_{\text{dec}}}\}$ (since $\Delta_k\le\Delta_{\max}$ and during a rejection the radius is multiplied by $\gamma_{\text{dec}}$) we obtain the constant $C$ stated in Theorem~\ref{thm:global}:
\[
C:=\frac{2}{c\,\eta_1(1-\eta_2)}\max\!\Bigl\{L,\;\frac{M}{\gamma_{\text{dec}}}\Bigr\}. \tag{17}
\]

\subsection*{Step 10 – Stationarity}
If the algorithm generated infinitely many iterates, the bound (16) forces the existence of a subsequence with $\|g_k\|\le\varepsilon$ for any $\varepsilon>0$. Hence $\liminf_{k\to\infty}\|g_k\|=0$. Because the radius never becomes negative and the model is continuously updated, the full limit holds:
\[
\lim_{k\to\infty}\|\nabla f(x_k)\|=0. \tag{18}
\]

All steps together establish Theorem~\ref{thm:global}. \qed

\section*{Rate Derivation (With Constants)}
From (16) we have the explicit iteration bound
\[
N(\varepsilon)\;\le\;\frac{\bar{B}}{\eta_1 c}\,\frac{f(x_0)-f^{\inf}}{\varepsilon^{2}}
\;=\;
\underbrace{\frac{M+\kappa\Delta_{\max}}{\eta_1 c}}_{\displaystyle C_1}
\;(f(x_0)-f^{\inf})\,\varepsilon^{-2}.
\]
If the model is exact ($\kappa=0$) and the Cauchy condition holds with $c=1$, the constant simplifies to $C_1=\frac{M}{\eta_1}$.

\section*{Special Cases}
\begin{enumerate}[label=(\alph*)]
\item \textbf{Convex quadratic $f(x)=\frac12x^{\top}Qx-b^{\top}x$ with $Q\succeq 0$.}  
The exact Newton step $p_k=-Q^{-1}\nabla f(x_k)$ always lies inside the trust region once $\Delta_k$ exceeds $\|p_k\|$, so the algorithm reduces to Newton’s method and converges linearly (or quadratically if $Q$ is positive definite and $\Delta_k$ is never restrictive).

\item \textbf{Strict saddle points.}  
If $x^{\star}$ is a saddle with a negative eigenvalue of $\nabla^{2}f(x^{\star})$, the model predicts a decrease along the direction of negative curvature. The trust‑region constraint caps the step length, but the ratio $\rho_k$ remains positive, guaranteeing that the algorithm moves away from the saddle in finitely many iterations.

\item \textbf{Gradient‑only implementation ($B_k=0$).}  
When $B_k$ is set to the zero matrix, the subproblem reduces to a scaled steepest‑descent step constrained by $\Delta_k$. The same analysis yields the classic $\mathcal{O}(\varepsilon^{-2})$ bound for gradient descent, showing that the trust‑region framework unifies first‑order and second‑order methods.
\end{enumerate}

\end{document}